% Lecture 10 - Page 7 (6-XII-1994).
% Fragment for % Lecture 10 - Page 7 (6-XII-1994).
% Fragment for % Lecture 10 - Page 7 (6-XII-1994).
% Fragment for % Lecture 10 - Page 7 (6-XII-1994).
% Fragment for \input{courses/pages/curs10-pagina7-en}.
% No preamble and no forced page break.
\bigskip

\subsection*{Lecture 10 - Page 7}

\noindent $aPa^{-1}=f_a(P)$, where $f_a$ is an automorphism.
\[
\Longrightarrow\quad aPa^{-1}\leq G
\quad\text{and}\quad |aPa^{-1}|=|P|.
\]

\bigskip
\noindent\fbox{\textbf{Sylow's Second Theorem}}. Under the same hypotheses
as Theorem 1, let $H$ be a $p$-subgroup of $G$, and let $P$ be a Sylow
$p$-subgroup of $G$.

\medskip
\noindent Then $(\exists)\,a\in G$ such that
\[
H\leq aPa^{-1}\in\operatorname{Syl}_p(G).
\]

\noindent If $P,Q\in\operatorname{Syl}_p(G)$, then
$(\exists)\,a\in G$ such that
\[
aQa^{-1}=P.
\]

\bigskip
\noindent\fbox{\textbf{Sylow's Third Theorem}}. Under the same hypotheses
as Theorem 1,
\[
|\operatorname{Syl}_p(G)|\mid |G|,
\]
\noindent and $|\operatorname{Syl}_p(G)|$ is of the form
$p^k+1$\footnote{The handwritten exponent appears to be $k$,
which gives $p^k+1$. The standard form of the theorem is $kp+1$.
The apparent manuscript reading has been retained pending verification.},
with
\[
|\operatorname{Syl}_p(G)|\equiv1\pmod p.
\]

\bigskip
\noindent\underline{Application of Sylow's theorems to the classification
of finite groups}.

\medskip
\noindent Given $n\in\mathbb{N}^{*}$, find a list of groups of order $n$,
$G_1,\ldots,G_q$, such that:
\begin{itemize}
\item for every group $G$ of order $n$, there is an $i$ such that $G\cong G_i$;
\item for every $i\neq j$, $G_i\not\cong G_j$.
\end{itemize}

\noindent $n=6:\quad (\mathbb{Z}_6,+,\hat 0),\quad(S_3,\circ,e).$

\bigskip
\noindent\fbox{\textbf{Theorem}}. If $G$ is a group such that
$|G|=n=pq$, $p>q$, $p,q$ are prime, and $q\nmid(p-1)$,
then
\[
G\cong\mathbb{Z}_{pq}.
\]

\medskip
\noindent For example, $n=35=7\cdot5$; then $G\cong\mathbb{Z}_{35}$.
% The example is heavily abbreviated in the manuscript.

\noindent Consider the case in which $q\mid(p-1)$.
\[
q=2,\qquad p=11.
\]

\bigskip
\noindent\underline{Lemma 2}. Let $G$ be a finite group and
$H$ a $p$-subgroup of $G$. Then
\[
[N_G(H):H]\equiv[G:H]\pmod p.
\]

\medskip
\noindent\underline{Proof of Lemma 2}: Take
\[
A=\{aH\mid a\in G\}.
\]
\noindent $H\times A\longrightarrow A$, an action by left translation:
\[
(h,aH)\longmapsto(ha)H.
\]
% The proof continues on the next manuscript page.
.
% No preamble and no forced page break.
\bigskip

\subsection*{Lecture 10 - Page 7}

\noindent $aPa^{-1}=f_a(P)$, where $f_a$ is an automorphism.
\[
\Longrightarrow\quad aPa^{-1}\leq G
\quad\text{and}\quad |aPa^{-1}|=|P|.
\]

\bigskip
\noindent\fbox{\textbf{Sylow's Second Theorem}}. Under the same hypotheses
as Theorem 1, let $H$ be a $p$-subgroup of $G$, and let $P$ be a Sylow
$p$-subgroup of $G$.

\medskip
\noindent Then $(\exists)\,a\in G$ such that
\[
H\leq aPa^{-1}\in\operatorname{Syl}_p(G).
\]

\noindent If $P,Q\in\operatorname{Syl}_p(G)$, then
$(\exists)\,a\in G$ such that
\[
aQa^{-1}=P.
\]

\bigskip
\noindent\fbox{\textbf{Sylow's Third Theorem}}. Under the same hypotheses
as Theorem 1,
\[
|\operatorname{Syl}_p(G)|\mid |G|,
\]
\noindent and $|\operatorname{Syl}_p(G)|$ is of the form
$p^k+1$\footnote{The handwritten exponent appears to be $k$,
which gives $p^k+1$. The standard form of the theorem is $kp+1$.
The apparent manuscript reading has been retained pending verification.},
with
\[
|\operatorname{Syl}_p(G)|\equiv1\pmod p.
\]

\bigskip
\noindent\underline{Application of Sylow's theorems to the classification
of finite groups}.

\medskip
\noindent Given $n\in\mathbb{N}^{*}$, find a list of groups of order $n$,
$G_1,\ldots,G_q$, such that:
\begin{itemize}
\item for every group $G$ of order $n$, there is an $i$ such that $G\cong G_i$;
\item for every $i\neq j$, $G_i\not\cong G_j$.
\end{itemize}

\noindent $n=6:\quad (\mathbb{Z}_6,+,\hat 0),\quad(S_3,\circ,e).$

\bigskip
\noindent\fbox{\textbf{Theorem}}. If $G$ is a group such that
$|G|=n=pq$, $p>q$, $p,q$ are prime, and $q\nmid(p-1)$,
then
\[
G\cong\mathbb{Z}_{pq}.
\]

\medskip
\noindent For example, $n=35=7\cdot5$; then $G\cong\mathbb{Z}_{35}$.
% The example is heavily abbreviated in the manuscript.

\noindent Consider the case in which $q\mid(p-1)$.
\[
q=2,\qquad p=11.
\]

\bigskip
\noindent\underline{Lemma 2}. Let $G$ be a finite group and
$H$ a $p$-subgroup of $G$. Then
\[
[N_G(H):H]\equiv[G:H]\pmod p.
\]

\medskip
\noindent\underline{Proof of Lemma 2}: Take
\[
A=\{aH\mid a\in G\}.
\]
\noindent $H\times A\longrightarrow A$, an action by left translation:
\[
(h,aH)\longmapsto(ha)H.
\]
% The proof continues on the next manuscript page.
.
% No preamble and no forced page break.
\bigskip

\subsection*{Lecture 10 - Page 7}

\noindent $aPa^{-1}=f_a(P)$, where $f_a$ is an automorphism.
\[
\Longrightarrow\quad aPa^{-1}\leq G
\quad\text{and}\quad |aPa^{-1}|=|P|.
\]

\bigskip
\noindent\fbox{\textbf{Sylow's Second Theorem}}. Under the same hypotheses
as Theorem 1, let $H$ be a $p$-subgroup of $G$, and let $P$ be a Sylow
$p$-subgroup of $G$.

\medskip
\noindent Then $(\exists)\,a\in G$ such that
\[
H\leq aPa^{-1}\in\operatorname{Syl}_p(G).
\]

\noindent If $P,Q\in\operatorname{Syl}_p(G)$, then
$(\exists)\,a\in G$ such that
\[
aQa^{-1}=P.
\]

\bigskip
\noindent\fbox{\textbf{Sylow's Third Theorem}}. Under the same hypotheses
as Theorem 1,
\[
|\operatorname{Syl}_p(G)|\mid |G|,
\]
\noindent and $|\operatorname{Syl}_p(G)|$ is of the form
$p^k+1$\footnote{The handwritten exponent appears to be $k$,
which gives $p^k+1$. The standard form of the theorem is $kp+1$.
The apparent manuscript reading has been retained pending verification.},
with
\[
|\operatorname{Syl}_p(G)|\equiv1\pmod p.
\]

\bigskip
\noindent\underline{Application of Sylow's theorems to the classification
of finite groups}.

\medskip
\noindent Given $n\in\mathbb{N}^{*}$, find a list of groups of order $n$,
$G_1,\ldots,G_q$, such that:
\begin{itemize}
\item for every group $G$ of order $n$, there is an $i$ such that $G\cong G_i$;
\item for every $i\neq j$, $G_i\not\cong G_j$.
\end{itemize}

\noindent $n=6:\quad (\mathbb{Z}_6,+,\hat 0),\quad(S_3,\circ,e).$

\bigskip
\noindent\fbox{\textbf{Theorem}}. If $G$ is a group such that
$|G|=n=pq$, $p>q$, $p,q$ are prime, and $q\nmid(p-1)$,
then
\[
G\cong\mathbb{Z}_{pq}.
\]

\medskip
\noindent For example, $n=35=7\cdot5$; then $G\cong\mathbb{Z}_{35}$.
% The example is heavily abbreviated in the manuscript.

\noindent Consider the case in which $q\mid(p-1)$.
\[
q=2,\qquad p=11.
\]

\bigskip
\noindent\underline{Lemma 2}. Let $G$ be a finite group and
$H$ a $p$-subgroup of $G$. Then
\[
[N_G(H):H]\equiv[G:H]\pmod p.
\]

\medskip
\noindent\underline{Proof of Lemma 2}: Take
\[
A=\{aH\mid a\in G\}.
\]
\noindent $H\times A\longrightarrow A$, an action by left translation:
\[
(h,aH)\longmapsto(ha)H.
\]
% The proof continues on the next manuscript page.
.
% No preamble and no forced page break.
\bigskip

\subsection*{Lecture 10 - Page 7}

\noindent $aPa^{-1}=f_a(P)$, where $f_a$ is an automorphism.
\[
\Longrightarrow\quad aPa^{-1}\leq G
\quad\text{and}\quad |aPa^{-1}|=|P|.
\]

\bigskip
\noindent\fbox{\textbf{Sylow's Second Theorem}}. Under the same hypotheses
as Theorem 1, let $H$ be a $p$-subgroup of $G$, and let $P$ be a Sylow
$p$-subgroup of $G$.

\medskip
\noindent Then $(\exists)\,a\in G$ such that
\[
H\leq aPa^{-1}\in\operatorname{Syl}_p(G).
\]

\noindent If $P,Q\in\operatorname{Syl}_p(G)$, then
$(\exists)\,a\in G$ such that
\[
aQa^{-1}=P.
\]

\bigskip
\noindent\fbox{\textbf{Sylow's Third Theorem}}. Under the same hypotheses
as Theorem 1,
\[
|\operatorname{Syl}_p(G)|\mid |G|,
\]
\noindent and $|\operatorname{Syl}_p(G)|$ is of the form
$p^k+1$\footnote{The handwritten exponent appears to be $k$,
which gives $p^k+1$. The standard form of the theorem is $kp+1$.
The apparent manuscript reading has been retained pending verification.},
with
\[
|\operatorname{Syl}_p(G)|\equiv1\pmod p.
\]

\bigskip
\noindent\underline{Application of Sylow's theorems to the classification
of finite groups}.

\medskip
\noindent Given $n\in\mathbb{N}^{*}$, find a list of groups of order $n$,
$G_1,\ldots,G_q$, such that:
\begin{itemize}
\item for every group $G$ of order $n$, there is an $i$ such that $G\cong G_i$;
\item for every $i\neq j$, $G_i\not\cong G_j$.
\end{itemize}

\noindent $n=6:\quad (\mathbb{Z}_6,+,\hat 0),\quad(S_3,\circ,e).$

\bigskip
\noindent\fbox{\textbf{Theorem}}. If $G$ is a group such that
$|G|=n=pq$, $p>q$, $p,q$ are prime, and $q\nmid(p-1)$,
then
\[
G\cong\mathbb{Z}_{pq}.
\]

\medskip
\noindent For example, $n=35=7\cdot5$; then $G\cong\mathbb{Z}_{35}$.
% The example is heavily abbreviated in the manuscript.

\noindent Consider the case in which $q\mid(p-1)$.
\[
q=2,\qquad p=11.
\]

\bigskip
\noindent\underline{Lemma 2}. Let $G$ be a finite group and
$H$ a $p$-subgroup of $G$. Then
\[
[N_G(H):H]\equiv[G:H]\pmod p.
\]

\medskip
\noindent\underline{Proof of Lemma 2}: Take
\[
A=\{aH\mid a\in G\}.
\]
\noindent $H\times A\longrightarrow A$, an action by left translation:
\[
(h,aH)\longmapsto(ha)H.
\]
% The proof continues on the next manuscript page.
